\documentclass[11pt]{article}
\usepackage[T1]{fontenc}
\usepackage[utf8]{inputenc}
\usepackage[spanish]{babel}
\usepackage{booktabs,hyperref,geometry}
\geometry{margin=2.5cm}
\title{El lenguaje T como optimización de la gestión del uso de la CPU durante el entrenamiento de redes neuronales artificiales}
\author{Guilherme Moura Fernandes\\\textit{Afiliación: [completar] -- Correo: [completar] -- ORCID: [si corresponde]}}
\date{Borrador científico privado -- no enviado}
\begin{document}\maketitle
\begin{abstract}
Presentamos la arquitectura auditable de Casulo/Lenguaje T y evaluamos, sin generalización universal, un ejecutor C especializado para una red neuronal superficial en CPU. La evidencia histórica congelada preservó parámetros, predicciones y pérdida respecto de PyTorch. En 21 pares históricos, la aceleración mediana fue 8,958 veces frente a PyTorch Eager (IC bootstrap pareado del 95\%: 8,075--9,872). El comparador Inductor disponible compiló solo la pasada hacia adelante y es secundario. No se midió energía eléctrica. El resultado apoya una hipótesis de reducción de overhead en un dominio restringido, no una superioridad general del lenguaje.
\end{abstract}
\section{Introducción} Motivación, problema, contribuciones verificables y preguntas.
\section{Trabajo relacionado} Compiladores tensoriales, especialización, fusión, benchmarking y energía.
\section{Arquitectura} Frontend general, IR, grafo, artefacto por contenido, backend CPU y ruta especializada separada.
\section{Metodología} Equivalencia primero, diseño pareado, calentamiento, sesiones, comparadores y pacto térmico.
\section{Resultados históricos} Medianas, intervalos, muestras crudas y equivalencia.
\section{Análisis causal y ablaciones} Especialización, Python, autograd, asignación, paralelismo y compilador.
\section{Amenazas a la validez} Un host, sesiones históricas, modelo sintético pequeño, float64, Inductor incompleto y sin julios.
\section{Conclusión} Resultado prometedor y específico; faltan múltiples hosts e integración del frontend.
\section*{Disponibilidad} Paquete privado verificable por hashes; publicación solo por decisión explícita del autor.
\bibliographystyle{plain}\bibliography{../references}
\end{document}

